\section{Bootstrap}
\label{sec:bootstrap}

Initializing a GFM actor requires discovering the capability lattice from
scratch. The lattice and the actor's ability to reason about $\volL(G)$ are
coupled: the actor cannot evaluate actions without a lattice, and the lattice
is populated through the actor's interactions. This section describes a
four-phase bootstrap procedure that uses both communication and observation
channels.

\subsection{Phase 1: Listen and Observe}

The actor begins with an empty lattice ($\C = \{\bot\}$) and low self-trust
$T_s$. It receives capability reports through the communication channel and
observes agent behavior through the observation channel. Each report or
observation that introduces a benchmarked capability adds a node to $\C$ with
weight $w(d) = \log_2(1 + \smax(d))$. All claims are initially weighted by
low trust $T_k$ (the cooling period from the foundational paper).

\subsection{Phase 2: Cross-Validate}

As reports accumulate, the actor compares the two channels. Agents whose
communicated capabilities match observed behavior build trust quickly
(consistent pattern from Section~\ref{sec:dual_channel}). Overclaiming agents
lose trust. Underclaiming agents trigger model expansion: the actor adds
observed capabilities that the agent did not report. The cross-validation phase
improves the quality of $\weff(d, k)$ and calibrates $T_k$ across the
population.

\subsection{Phase 3: Discover Structure}

The actor discovers subsumption relationships through two mechanisms:

\emph{Logical analysis.} If capability descriptions in $L$ have compositional
structure, the actor can infer subsumption from the descriptions themselves.
``Can run a marathon in under 4 hours'' logically implies ``can run a mile.''
The shared language $L$ must be expressive enough for this inference; formal
specification languages or capability ontologies support it directly.

\emph{Empirical observation.} If every agent observed to have capability $d_1$
also demonstrates $d_2$, the actor infers $d_2 \leq d_1$ (with confidence
proportional to the number of confirming observations and the diversity of
agents observed). This is inductive: a sufficiently large and diverse
observation set makes the inference reliable.

As subsumption edges are added, the actor recomputes independent weights
$\iw(d)$ via full M\"{o}bius inversion ($O(|\C|^3)$ for general lattices,
$O(|\C| \cdot k)$ for trees; Proposition~\ref{prop:tractability}).

\subsection{Phase 4: Expand}

With a populated lattice and calibrated trust, the actor begins proposing
actions and evaluating them via $\Deltahat(\pi)$
(Equation~\eqref{eq:weighted_estimator}). Self-trust $T_s$ increases as the
actor's predictions ($\Deltahat$) match observed outcomes ($\Delta \volL$
computed post-action). Both channels continue refining the lattice throughout
the actor's operation: new capabilities are discovered, benchmark scores are
updated, and subsumption relationships are confirmed or revised.

\subsection{Convergence}

\begin{proposition}[Lattice Growth]
\label{prop:bootstrap_convergence}
Under assumptions of positive communication frequency and non-zero observation
    coverage:
\begin{enumerate}
\item[\textbf{(a)}] The lattice node set grows monotonically: $\C(t) \subseteq
    \C(t+1)$ for all $t$.
\item[\textbf{(b)}] Adding new capabilities to $G$ (without changing
    subsumption structure) is $\volL$-non-decreasing.
\item[\textbf{(c)}] Discovering subsumption edges can \emph{decrease}
    $\volL(G)$: if $d_1$ and $d_2$ were treated as unrelated and the edge
    $d_1 \leq d_2$ is later discovered, the downward closure $\dc G$ does not
    change but the M\"{o}bius-inverted weights redistribute, and $\volL(G)$
    may drop from $w(d_1) + w(d_2)$ to $w(d_2)$.
\item[\textbf{(d)}] Updating a benchmark to a finer scale ($\smax' > \smax$)
    is $\volL$-non-decreasing: $\volL$ increases if $d$ is maximal in $\dc G$
    and is unchanged otherwise.
\end{enumerate}
\end{proposition}

\begin{proof}
(a) Nodes and edges are never removed from $\C$.

(b) If $G(t+1) = G(t) \cup \{d\}$ for a new capability $d$ with no new
subsumption edges, then $\dc G(t+1) \supseteq \dc G(t)$ and all independent
weights are unchanged. By M3, $\volL(G(t+1)) \geq \volL(G(t))$.

(c) Counterexample: let $d_1$ = ``write a function'' ($w = 6.66$) and $d_2$ =
``build a web app'' ($w = 9.97$). Before discovering $d_1 \leq d_2$, the
capabilities are lattice-disjoint: $\volL(\{d_1, d_2\}) = 6.66 + 9.97 =
16.63$. After adding the edge, M\"{o}bius inversion gives $\iw(d_1) = 6.66$,
$\iw(d_2) = 9.97 - 6.66 = 3.31$, and $\volL(\{d_1, d_2\}) = 6.66 + 3.31 =
9.97 = w(d_2)$. The volume decreased by $6.66$. This is epistemically correct:
the actor previously overcounted because it did not know that $d_1$'s
measurement content was already captured by $d_2$. Structural refinement
corrects the overcount.

(d) Increasing $\smax(d)$ increases $w(d) = \log_2(1 + \smax(d))$ and therefore
increases $\iw(d)$. However, for any ancestor $d' > d$ in $\dc G$, the
M\"{o}bius sum for $\iw(d')$ includes $w(d)$ with coefficient $\mu(d, d')$,
which may be negative, decreasing $\iw(d')$. The net effect on $\volL(G)$
depends on whether $d$ is maximal in $\dc G$: if $d$ has no descendant in
$\dc G$ above it, $\volL$ increases by $\Delta w(d)$; if $d$ is non-maximal,
the increase in $\iw(d)$ is absorbed by offsetting decreases in ancestors'
independent weights, and $\volL(G)$ is unchanged
(Appendix~\ref{app:proof_bootstrap}).
\end{proof}

\paragraph{Interpretation.} Part (b) shows that discovering new capabilities
always increases $\volL$. Part (d) shows that refining benchmarks increases
$\volL$ when the refined capability is maximal in the closure, and leaves it
unchanged otherwise (the increase is absorbed by ancestors). Part (c) shows
that discovering structure (subsumption) can decrease $\volL$ by correcting
overcounting. The net trajectory of $\volL$ during bootstrap depends on the
balance between capability discovery (expanding) and structural refinement
(potentially contracting). In practice, early bootstrap is dominated by (b)
(rapid discovery of new capabilities), while later phases are dominated by (c)
(structural refinement that tightens the measure). The long-run $\volL$
converges toward the correct volume given the actor's full structural
knowledge, which may be lower than the naive estimate from the early
structureless phase.

\paragraph{Binary default and deferred structure.} Every capability that
enters the lattice through communication or observation immediately receives
the binary default weight $w(d) = 1$ bit ($\smax = 1$): the actor knows the
capability exists. The capability enters as an \emph{unstructured atom} with no
subsumption edges, contributing $\iw(d) = w(d) = 1$ to $\volL(G)$ immediately.

Subsumption edges are added in a separate step (Phase~3) only after the
non-negativity check passes: for a proposed edge set making $d$ a multi-cover
node, the actor verifies $\iw(d) \geq 0$ under M\"{o}bius inversion before
committing the edges. If the check fails (e.g., a diamond node with two
binary-default prerequisites gives $\iw = 1 - 1 - 1 < 0$), the edges are
deferred until the capability's benchmark is refined to a level where the check
passes. The capability continues contributing as a lattice-disjoint atom in the
interim.

This sequencing ensures two properties simultaneously: (1) every discovered
capability contributes to $\volL(G)$ from the moment of discovery, and (2)
$\iw(d) \geq 0$ is maintained for all nodes at all times. Developing a finer
benchmark ($\smax > 1$) increases $w(d)$ beyond the binary default, eventually
enabling deferred edges to be added. The remaining zero-weight gap consists
only of capabilities the actor has not yet learned about through either
channel.
